\documentclass[11pt]{article}

\usepackage[margin=1in]{geometry}
\usepackage{amsmath}
\usepackage{amssymb}
\usepackage{booktabs}
\usepackage{hyperref}

\title{Notes on Scaling the Quasi-Symmetry DMRG Workflow\\
to Larger Orbital Spaces}
\author{}
\date{July 2026}

\begin{document}
\maketitle

\section{Purpose}

The current workflow is practical for the full STO-3G calculations considered
so far, whose largest examples contain ten spatial orbitals and twenty qubits.
It already avoids several unnecessary costs: DMRG is performed entirely with
\texttt{pyblock2}, CISD determinants are added to the initial MPS in bounded
batches, Clifford circuits and qubit permutations act directly on the MPS,
and each Hamiltonian representation is processed and saved separately.

Larger basis sets introduce qualitatively different bottlenecks.  The main
issue is not DMRG alone.  Full-CI reference states, entropy calculations,
CISD-to-MPS conversion, symmetry scoring, transformed Pauli Hamiltonians, and
MPO construction can each become the limiting step.  This note records a
possible sequence of changes rather than prescribing a final implementation.

\section{Remove full CI from the required workflow}

The size of the full-CI space grows combinatorially with the numbers of
electrons and orbitals.  Full CI should therefore become an optional
small-system validation method rather than a required input.

For larger systems, a tightly converged spin-adapted fermionic DMRG
calculation can provide both the reference energy and reference MPS.  The
reference energy used in the chemical-accuracy test would then be labeled by
its source, for example FCI, selected CI, variational DMRG, or extrapolated
DMRG.  The reported criterion remains
\begin{equation}
  \left|E(\chi)-E_{\mathrm{ref}}\right|\leq 1.6\ \mathrm{mHa},
\end{equation}
but it no longer assumes that $E_{\mathrm{ref}}$ is available from FCI.

The software interface should consequently accept a general reference-state
object.  For small systems this may be a sparse CI vector; for larger systems
it should normally be a saved MPS together with its energy, quantum numbers,
and convergence information.

\section{Calculate entropies from an MPS}

The current preprocessing calculates entropies and mutual information from a
saved FCI state.  Instead, one- and two-site reduced density matrices can be
contracted directly from a converged reference MPS.  This gives the one-site
entropies and pairwise mutual information needed for Fiedler ordering without
constructing a vector of length $2^{n_q}$.

For a transformed representation, the same unitary circuit used for the
Hamiltonian can be applied to a copy of the reference MPS.  Entropies can then
be measured from the transformed MPS.  This procedure applies equally to
Clifford circuits, qubit permutations, and orbital rotations decomposed into
Givens rotations.  Every entropy result should record the MPS bond dimension
and discarded weight used in this preprocessing step.

\section{Make the orbital space explicit}

In a larger basis, retaining every orbital will often be neither necessary
nor affordable.  Molecular input and saved metadata should explicitly record
the basis, frozen orbitals, active orbitals, active electron number, target
spin, and original orbital ordering.  A useful interface should support both
the full orbital space and a selected active space without changing the
subsequent symmetry and DMRG code.

The active-space definition must also be carried through the CISD state,
fermionic MPO, Jordan--Wigner Hamiltonian, and reference energy.  Silent
differences between these spaces would invalidate energy comparisons.

\section{Improve CISD-to-MPS conversion}

Although CISD is polynomial in system size, its number of double excitations
scales approximately as
\begin{equation}
  O\!\left(n_{\mathrm{occ}}^2 n_{\mathrm{virt}}^2\right).
\end{equation}
The present method retains every CISD determinant, sorts them by coefficient
magnitude, adds them sequentially to an MPS, and compresses the growing sum in
batches.  This is reliable for the current systems but may become slow for
large virtual spaces.

Possible improvements include balanced-tree addition of determinant MPSs,
adaptive batch sizes, an intermediate bond-dimension cap, and construction by
applying single- and double-excitation operators to the Hartree--Fock MPS.
Disk-backed intermediate tensors may also be useful.  Selected-CI or
coefficient truncation can be offered as an explicit optional approximation,
but the retained norm, discarded norm, and change in energy must be reported.

\section{Avoid large intermediate Pauli operators}

An electronic Hamiltonian contains $O(n_{\mathrm{orb}}^4)$ integral terms.
Jordan--Wigner mapping and Clifford conjugation can make the explicit Pauli
representation expensive even when the corresponding fermionic MPO remains
manageable.  Materializing the complete transformed \texttt{QubitOperator}
and all intermediate MPO sums may therefore become the dominant memory cost.

The transformation and MPO builder should eventually operate as a streaming
pipeline: transform terms in chunks, combine identical Pauli words
immediately, and release each chunk after it has been incorporated into the
MPO.  Coefficient screening should be optional and controlled by an explicit
error tolerance.  Density fitting, Cholesky decomposition, and other
factorized representations of the two-electron interaction should also be
investigated, especially if the factorized structure can be preserved during
MPO construction.

For sufficiently large orbital spaces, a Pauli-space DMRG calculation may
remain much more expensive than the spin-adapted fermionic calculation.  That
difference is scientifically relevant and should not be hidden by using
different uncontrolled truncations in the two representations.

\section{Scale the symmetry search}

The present beam-search safeguards---50 pairwise seed terms and a candidate
pool capped at 1000 operators---limit the combinatorial part of the search.
The remaining cost comes from repeatedly scoring candidates against a large
Hamiltonian and state.

Useful improvements include caching Pauli commutation relations in binary
symplectic form, evaluating candidate scores incrementally, screening small
Hamiltonian terms during pool generation, and parallelizing independent
candidate evaluations.  The initial candidate pool, its ranking scores, and
the state of a partially completed search should be saved so that a long
search can be resumed exactly.

When the Hamiltonian and state are transformed by a known unitary, objective
values can sometimes be evaluated in the original frame by transforming only
the candidate Pauli operator back to that frame.  For Clifford transformations
this maps a Pauli product to another Pauli product.  It can avoid constructing
a large sparse matrix for every transformed Hamiltonian and is especially
useful for repeated symmetry searches.

\section{Control the cost of applying unitary circuits}

Applying a long circuit to an MPS can increase its bond dimension before DMRG
begins.  The circuit layer should report the discarded singular-value weight,
the maximum intermediate bond dimension, and the energy error introduced by
compression.  Separate bond caps may be appropriate for reference-state
transformation and for the later DMRG calculation.

Additional improvements include reordering commuting gates, reducing swap
routing, applying independent gates in layers, and variationally compressing
the state after the complete circuit.  These operations should use the common
unitary interface so that Clifford transformations and orbital rotations
follow the same execution and validation path.

\section{Exploit exact symmetries where available}

The fermionic calculations already use particle-number and spin-adapted
$\mathrm{SU}(2)$ symmetry.  The general Pauli calculations currently do not
enforce particle number or total spin.  For larger systems, exact mapped
symmetries could be used to restrict the qubit calculation or eliminate fixed
qubits.  Approximate symmetries cannot be fixed without introducing an
additional approximation, so exact and approximate generators must remain
clearly distinguished.

\section{Suggested implementation order}

The following order removes the most restrictive bottlenecks first:
\begin{enumerate}
  \item introduce a reference-state interface supporting saved MPSs;
  \item compute entropies and mutual information directly from an MPS;
  \item generalize molecular input and active-space metadata;
  \item improve and instrument CISD-to-MPS conversion;
  \item add streaming transformed-Hamiltonian and MPO construction;
  \item cache and checkpoint the symmetry search; and
  \item optimize circuit routing and MPS compression.
\end{enumerate}

This sequence preserves the current small-system workflow as a validation
case while allowing the reference calculation, entropy analysis, and warm
start to be replaced independently as the orbital space grows.

\end{document}
